% PVLDB Vol. 20 (VLDB 2027) submission. Template: vldbproceedings/VLDB-Template
% (acmart.cls v2.19, pinned; see TEMPLATE_VERSION). Do not edit acmart.cls or pvldb.sty.
%
% Research track, regular research paper: up to 12 pages excluding references.
% Rolling deadline on the 1st of each month (abstract by the 25th of the previous
% month), final deadline 2027-03-01. Single-blind: names and affiliations required.

\documentclass[sigconf, nonacm]{acmart}

%%% do not modify the following VLDB block %%
%%% VLDB block start %%%
\usepackage{pvldb}
%%% VLDB block end %%%

\renewcommand\vldbdoi{XX.XX/XXX.XX}
\renewcommand\vldbpages{XXX-XXX}
\renewcommand\vldbavailabilityurl{https://github.com/dtenwolde/duckdb-parser-experiments}

\usepackage{booktabs}
\usepackage{balance}
\usepackage{xcolor}
\usepackage{listings}

% Draft helpers. Remove before submission.
\newcommand{\todo}[1]{\textcolor{red}{[TODO: #1]}}
\newcommand{\num}[1]{\textcolor{blue}{#1}} % numbers that must be re-measured on the final commit

\lstdefinelanguage{peg}{
  morekeywords={},
  sensitive=true,
  morecomment=[l]{\#},
  morestring=[b]',
}
\lstset{basicstyle=\ttfamily\small, columns=fullflexible, keepspaces=true, breaklines=true}

\begin{document}
\title{Runtime-Extensible SQL Parsing: Design, Use Cases, and Lessons from a Production PEG Parser}
% Working title. Alternatives: "A Runtime-Extensible Parser for a Production Database System".

\author{Daniel ten Wolde}
\affiliation{%
  \institution{CWI}
  \city{Amsterdam}
  \country{The Netherlands}
}
\email{daniel.ten.wolde@cwi.nl}

% \todo{co-authors and affiliations: confirm with supervisor and DuckDB Labs}

\begin{abstract}
\todo{Abstract. Claim: a runtime-extensible PEG parser replaces the compile-time
yacc parser of a production DBMS, lets loadable extensions add syntax (SQL/PGQ) or
whole dialects (SaneQL) without core changes, and does so with measured cost.
No citations in the abstract. Must match the CMT abstract exactly.}
\end{abstract}

\maketitle

%%% do not modify the following VLDB block %%
%%% VLDB block start %%%
\vldbtopmatter
%%% VLDB block end %%%

\section{Introduction}
\label{sec:intro}

\todo{Motivation: SQL front ends are compile-time artifacts (yacc/bison); every new
feature needs new syntax (SQL:2023 has 409 keywords~\cite{neumann2024saneql});
extensions cannot add syntax without forking the grammar. The CIDR 2025 vision
paper~\cite{muhleisen2025runtime} argued for runtime-extensible parsers. This paper
reports on building one and shipping it as the default parser of DuckDB.}

\todo{Contributions list:
(1) design of a runtime-extensible PEG parser for a production DBMS:
grammar changes and dialects applied at load time, heap-based execution
for embedded hosts, differential testing against the legacy parser until parity;
(2) three use cases built on it: SQL/PGQ via DuckPGQ (GrammarChange API),
SaneQL as a loadable dialect (DialectExtension API), and \todo{third: agentic
extension (INSPECT / WITH OUTPUT) or another DialectExtension user};
(3) a first version of PEG syntax-error diagnostics evaluated on a mutation corpus;
(4) performance engineering techniques and their measured impact, including
heap-based versus stack-based execution;
(5) lessons learned from \num{~2{,}100} parser commits.}

\section{Background}
\label{sec:background}

\subsection{SQL Parsing in Database Systems}
\todo{libpg\_query / bison in DuckDB before 2025; Calcite~\cite{begoli2018calcite},
Spark SQL~\cite{armbrust2015sparksql} parser sections; nobody reports parse share
of end-to-end latency~\cite{kersten2021tidy}.}

\subsection{Parsing Expression Grammars}
\todo{PEG~\cite{ford2004peg}, packrat~\cite{ford2002packrat}, ordered choice,
no ambiguity by construction, prefix capture. Extensibility claims in
Rats!~\cite{grimm2006rats}; FIRST-set choice in LPeg~\cite{medeiros2008lpeg};
none evaluated on a SQL-scale grammar in a production system.}

\subsection{The Runtime-Extensible Parser Vision}
\todo{Summarise~\cite{muhleisen2025runtime}; what this paper adds beyond it.}

\section{A Runtime-Extensible PEG Parser}
\label{sec:design}

\subsection{Grammar Representation}
\todo{\texttt{.gram} files, \num{1{,}107} rules, keyword categories as bit flags,
terminal rules vs non-terminal rules, macros (Parens/List).}

\subsection{Matcher Engine}
\todo{Matcher graph built from the grammar at load time; MatchStack; explicit
heap-allocated frames instead of recursion. Why: embedded hosts (Python, R, WASM)
with small or unknown stack sizes; deep nesting; the parser initially was
recursive (stack-based) and was rewritten. Memoization (packrat) must replay
side effects (token annotations) or caches diverge. Rollback boundary, not grammar
node, is the unit of shareable state.}

\subsection{Transformer}
\todo{Parse tree to DuckDB AST; generated trampolines (\num{23{,}702} lines, generated)
vs hand-written transforms (\num{9{,}778} lines); PEGRuleId.}

\subsection{Extension Interfaces}
\todo{GrammarChange API (AddRule, ReplaceRule, AddChoice, PrependChoice,
terminal overrides) used by DuckPGQ; DialectExtension API
(ApplyGrammarChanges, CreateTokenizer, CreateMatcherFactory) used by SaneQL;
\texttt{SET current\_dialect}, \texttt{SET active\_grammar\_extensions};
statement-scoped grammar changes as open problem (SaneQL \texttt{defun}).}

\subsection{Equivalence Gate}
\todo{Differential testing against libpg\_query: \num{90k} statement fingerprints +
\num{10k} synthetic matcher graphs; how it caught ordered-choice ordering bugs that
yacc would have reported as conflicts.}

\section{Use Cases}
\label{sec:usecases}

\subsection{SQL/PGQ with DuckPGQ}
\todo{Before: fork of DuckDB with patched libpg\_query grammar
(cwida/duckdb-pgq)~\cite{tenwolde2023duckpgq}. After: GrammarChange API, loadable
against upstream DuckDB. Lines of code before/after, maintenance cost (point
release merges), what SQL/PGQ syntax needed (fused operator tokens, label
expressions, quantifiers). Pitfall: wrapper rules without transform process.}

\subsection{SaneQL as a Loadable Dialect}
\todo{SaneQL~\cite{neumann2024saneql} as DialectExtension: \num{24} PEG rules,
\num{753} C++ lines, no core changes, all 22 TPC-H queries pass. Table from
saneql/PAPER\_NOTES.md: reference translator 5{,}028 lines (fair subset
grammar+lexer+semana 2{,}641) vs 753. Framing: syntax is a thin replaceable layer
when the semantic layer is shared. Report all caveats: \texttt{defun} not
implemented (closed operator keyword set), host SQL lexical decisions leak into
the dialect (reserved words, \texttt{||}, \texttt{a.desc()}), \texttt{let} as three
rewrites, API gaps (AddTerminalRuleOverride, undocumented parse-tree shape).}

\subsection{\todo{Third Use Case}}
\todo{Options: (a) agentic SQL extensions as parse-time rewrites via GrammarChange:
\texttt{INSPECT} (schema+stats+samples) and a \texttt{WITH OUTPUT} clause
(max\_rows, truncate, total count); (b) another DialectExtension user
(Spark dialect is private; GGSQL possible collaborator; pipe syntax~\cite{shute2024pipe}
as a small dialect would be a clean third point). Decide by what can be built and
measured in time.}

\section{Syntax Error Diagnostics}
\label{sec:errors}

\todo{First version. PEGs report the farthest failure position; libpg\_query
reports the yacc error token. Our approach: reuse the autocomplete suggestion walk
to compute expected tokens at the failure point; verified "Did you mean" keyword
corrections (one-token look-back for typos swallowed as aliases); verified
single-token "Missing X" insertions; rule labels ("an expression", "a type").
Design constraints: the suggestion pass must keep the packrat cache on or it goes
exponential on unclosed nesting (\num{56 s} at depth 20); the re-parse must use
BUILD\_PARSE\_RESULT mode (RECOGNIZE\_ONLY skips first-set pruning, \num{5x} work);
one extra parse per error is the floor. Agent-mode JSON errors. Relate to
PEG error reporting~\cite{maidl2016error} and recovery~\cite{medeiros2018recovery}.
Evaluation in \autoref{sec:eval-errors}.}

\section{Performance Engineering}
\label{sec:performance}

\todo{Performance is still improving; this section describes techniques and the
evaluation reports their impact on whatever commit the final run uses. Techniques:
FIRST sets, position-indexed memo table, arena allocation, cursor sharing
(runtime left-factoring), generated transform schedules, second-token rejection,
forwarding rules, per-step inlining. Heap-based frames vs the original recursive
implementation: cost and benefit. Negative results worth reporting: frame counts do
not predict time (lookahead cut steps 39\% but time 2.5\%); smaller structs were
slower; per-step cost (~55 ns) in engine bookkeeping dominates, only ~4\% of matcher
self time compares tokens; PEG allocates ~390 KB/query vs 27 KB for libpg\_query.
Open direction: eliminate the intermediate capture tree with rollback-safe
semantic actions.}

\section{Evaluation}
\label{sec:eval}

\todo{Setup: hardware, DuckDB commit (rolling; final numbers re-measured on one
commit), libpg\_query baseline (v1.5.5 binary or legacy path in the same binary),
interleaved runs, >=3 rounds, medians. Artifacts in this repository.}

\subsection{Extensibility}
\label{sec:eval-ext}
\todo{Table: use case, rules added/replaced, C++ lines, core changes (0), tests
passing. DuckPGQ before/after. SaneQL vs reference translator. Grammar construction
time at load (matters because extensions rebuild the matcher graph).}

\subsection{End-to-End Parse Share}
\label{sec:eval-share}
\todo{parser\_ms / latency\_ms on TPC-H, TPC-DS at SF 0, 0.1, 1 (scripts/benchmark\_parser.sh),
ClickBench, and short queries (point lookups, where parse share is largest).
PEG vs libpg\_query per query; heatmaps from analysis/.}

\subsection{Parser Microbenchmarks}
\label{sec:eval-micro}
\todo{Phase breakdown per token: tokenizer (PEG 45 ns vs flex 65 ns), matcher
(412 vs 27 ns), transformer (142 vs 44 ns). Ablation: impact of each technique
measured by cherry-picking commits in branch order (cross-revision harness);
heap-based vs stack-based execution.}

\subsection{Syntax Error Quality}
\label{sec:eval-errors}
\todo{Mutation corpus: ~37k statements from the test suite, TPC-H and TPC-DS;
score error position accuracy and hint hit rate for PEG, libpg\_query (v1.5.5)
and PostgreSQL 17 (pglast). Report common-only and per-parser accepted subsets.
Caveat: duplicated opening bracket is reported late by every parser.}

\subsection{Bug-Fix Taxonomy}
\label{sec:eval-bugs}
\todo{~453 fix-like of ~2{,}134 parser commits since April 2025 (filter merge/format
noise). Hypothesis: ordered-choice ordering bugs replace yacc conflict diagnostics.}

\section{Lessons Learned}
\label{sec:lessons}

\todo{A shared semantic layer imports the host language's lexical decisions.
Terminal-level overrides are a distinct extension point from rule-level ones.
Memoization must replay effects. Judge "reduce attempts" ideas by time, not frame
counts. Ordered choice moves grammar bugs from compile-time conflicts to runtime
ordering; differential testing against a reference parser is mandatory. Statement-scoped grammar changes are
needed for user-defined syntax (\texttt{defun}). Parse-tree shape must be documented
as API.}

\section{Related Work}
\label{sec:related}

\todo{SQL parsing at DB venues: only~\cite{muhleisen2025runtime} and
SQLFlex~\cite{an2026sqlflex} (LLM segmentation for dialect robustness; the
alternative to runtime-extensible grammars). Dialect translation with LLMs
(CrackSQL, DialectEvo, Dial) as application pull. Language design: pipe
syntax~\cite{shute2024pipe}, SaneQL~\cite{neumann2024saneql}. Parsing as a DB
performance topic: Mison~\cite{li2017mison}, simdjson~\cite{langdale2019simdjson}.
PEG literature: \cite{ford2004peg,ford2002packrat,grimm2006rats,medeiros2008lpeg,
mizushima2010cut,maidl2016error,medeiros2018recovery}. Extension case study:
DuckPGQ~\cite{tenwolde2023duckpgq}.}

\section{Conclusion}
\label{sec:conclusion}

\todo{Summary of claims and evidence; future work: statement-scoped grammar
changes, capture-tree elimination, error recovery.}


\begin{acks}
\todo{Funding and acknowledgements.}
\end{acks}

\balance
\bibliographystyle{ACM-Reference-Format}
\bibliography{references}

\end{document}
